\scanpage{11}\mspage{10}%
ci inspirent.

\uline{Axiomes}\quad On peut prendre pour \unsure{ens} de base \unsure{ds} une axiomatique des ateliers, soit $\unsure{\mathcal{F}|\mathscr{F}}$, soit $M$, soit $\mathcal{L}$. On trouve des données les plus simples sur $\mathcal{F}$, les \deleted{\unsure{plus}} \unsure{sophistiquées|sophistiqués} sur $\mathcal{L}$. Par contre, les axiomes essentiels sont les moins transparents quand on \unsure{reste|entre} \unsure{plus ou moins} dans $\mathcal{F}$, \unsure{le|les} plus \unsure{transparent|transparents} quand on travaille dans $\mathcal{L}$ (ou peut-être pareil pour $M$ ?).

\uline{Sur $\mathcal{F}$}, on peut prendre comme structure de base, au choix
\[
\left\{\begin{aligned}
 &\text{a)}\ \leq,\ \ll\\
 &\text{b)}\ \leq,\ \preccurlyeq \quad \longleftarrow\\
 &\text{c)}\ \ll,\ \deleted{\unsure{F|\mathcal F}}\quad
 \underbrace{\text{et}\quad\mathcal{F}\longrightarrow\mathrm{P}(\mathcal{F})}_{
 \begin{smallmatrix}
 \text{i.e. la \uline{relation}}\\[-1pt]
 \text{\uline{d'incidence}}\\
 \text{``$F$ est une strate de $G$''}\\
 \text{(i.e. $F\in M$ \uline{et} $F\leq G$)}
 \end{smallmatrix}}
\end{aligned}\right.
\quad
\begin{aligned}
 &\text{(notée }F\longmapsto\widetilde F,\ \text{le}\\
 &\text{\unsure{déployement|déploiement} de }F\\
 &\text{vu comme partie de }\mathcal{F}\text{)}\\
 &\text{\deleted{i.e. la \unsure{description|désignation} \inserted{relation \unsure{de type|d'inc.}}}}
\end{aligned}
\tag*{(1.16)}
\]
\margin{(20. juin) Non, \unsure{par mon nouveau point de vue} ça n'implique plus $\ll$}

\deleted{\unsure{Je}} Avant de préciser ces points, je vais grouper dans un diagramme les six relations fondamentales dans $\mathcal{F}$, \unsure{on écrira|en écrivant}
\[
F\triangleleft G\quad
\underset{\unsure{\mathrm{par}}}{\overset{\unsure{\mathrm{dfn}|\mathrm{déf}}}{\Longleftrightarrow}}\quad
\text{``$F$ incident à $G$''}\quad\text{i.e.}\quad F\leq G\ \text{et}\ F\in M
\tag*{(1.17)}
\]
